\chapter{Justificación}\label{chapter:PA}

Abordar el diseño de sensores plasmónicos desde el dominio de la Inteligencia Artificial (IA) representa una necesidad oportuna por varias razones críticas estipuladas en el cruce de la fotónica y la ciencia de datos. En primer lugar, la eficiencia de predicción sin precedentes aportada por las arquitecturas de aprendizaje profundo (Deep Learning) permitirá reemplazar el modelo principal físico clásico de base matricial (TMM) con un \textit{Modelo Sustituto} (Surrogate Model), reduciendo el tiempo de inferencia de minutos usando cálculo numérico estándar CPU, a escasos milisegundos de propagación directa en inferencia acelerada mediante GPUs como la NVIDIA RTX 4060 Ti con núcleos Tensor \cite{surrogate_optics}.

La segunda contribución innovadora de peso en este estudio es aplicar el modelado computacional para el \textit{Diseño Inverso} (Inverse Design). A través de modelos secuenciales, el sistema IA resultante servirá como un oráculo de ingeniería capaz de predecir la estructura física necesaria del sensor de plasmón basándose en las curvas de sensibilidad objetivo, omitiendo así las exhaustivas heurísticas de ensayo y error convencionales \cite{ml_plasmonic_inverse}.

En tercer lugar, el marco de implementación contemplado incluye el despliegue del estado del arte de las Redes Neuronales Informadas por la Física (Physics-Informed Neural Networks - PINNs) y de las metodologías de IA Explicable (XAI). Las PINNs limitan las aproximaciones a configuraciones materialmente y físicamente posibles asegurando el respeto estricto de las Ecuaciones de Maxwell disminuyendo masivamente el volumen de datos de entrenamiento empíricos requeridos \cite{pinn_nanophotonics}. En paralelo, la adopción de análisis de explicabilidad usando herramientas post-hoc como SHAP (SHapley Additive exPlanations), abrirán la caja negra algorítmica y dotarán al experto de grafos interpretables para descubrir exactamente qué leyes fotónicas o de espesor están incidiendo a una mejor sensibilidad \cite{xai_spr_sensors}.

Finalmente, a nivel metodológico la estructuración y apertura del repositorio software agrupará estos pilares en un único simulador de resonancia SPR de extremo a extremo. Dicha aproximación integrativa en IA aplicada enriquecerá drásticamente el panorama investigativo, ofreciendo una plataforma para que futuros laboratorios y grupos de ingenieros aborden retos físicos similares en otras esferas tecnológicas mediante aprendizaje maquinal.
